\documentclass[11pt,a4paper]{ctexart}
\usepackage[margin=2.4cm]{geometry}
\usepackage{amsmath,amssymb}
\usepackage{fancyhdr}
\usepackage{enumitem}
\linespread{1.18}
\setlength{\headheight}{15pt}
\setlength{\parskip}{0pt}

\pagestyle{fancy}
\fancyhf{}
\fancyhead[L]{\small Team 2303950}
\fancyhead[R]{\small \thepage}
\renewcommand{\headrulewidth}{0.4pt}
\fancypagestyle{plain}{\fancyhf{}\renewcommand{\headrulewidth}{0pt}}

\ctexset{
  section/format = \Large\bfseries\raggedright,
  subsection/format = \large\bfseries\raggedright,
  subsubsection/format = \normalsize\bfseries\raggedright
}
\setcounter{secnumdepth}{3}

\begin{document}

%====================== 摘要页 ======================
\thispagestyle{empty}
\begin{center}
{\small 2023 MCM/ICM\quad 选题 A\quad 队伍控制号 2303950\quad 摘要页（Summary Sheet）}\\[18pt]
{\LARGE\bfseries 干旱胁迫植物群落的模拟方法}\\[8pt]
{\large （A Methodology to Simulate Drought-Stricken Plant Communities）}\\[10pt]
\end{center}

\begin{center}
{\large\bfseries 摘\quad 要}
\end{center}

\vspace{6pt}
干旱地区植物群落的研究是生态环境保护中不可或缺的组成部分。在这类区域中，一年生植物凭借其独特的性状而成为植被的重要组成部分。本文提出了一套针对一年生植物生长、竞争与促进作用的模拟方法，该方法可推广应用于多种情形。

首先，我们建立了一年植物的生长模型。我们先确定 Logistic 模型中土壤含水量与自然增长率之间的关系；随后基于随机降雨模型建立土壤含水量的微分方程；最后给出一年生植物的繁殖规律。为估计随机降雨模型中的统计参数，我们依据最小二乘法的正规方程（Normal Equation）求得最优参数取值。

在生长模型的基础上，我们建立了多物种的竞争模型，以模拟植物之间的根系竞争。对于 Lotka-Volterra 模型中的竞争系数，我们创新性地确定了它与植物根系生长、水分利用能力之间的关系。结合根系的结构特征，我们对模型中的参数作出了合理估计。

在竞争模型的基础上，我们进一步建立了灌木植物对草本植物的促进作用模型。考虑到灌木植物的提水作用（water-lifting），我们对竞争模型中的吸水速率进行了修正。至此，我们的模型已得到完整阐述，模型中的参数也已得到恰当的估计。

此后，我们采用四阶 Runge-Kutta 算法对微分方程进行数值求解。问题要求中的各项问题已得到如下解答：

\begin{itemize}
\item 基于我们 TOPSIS 模型的结果，至少需要八个植物物种才能为群落提供整体效益。随着植物物种数目的增加，群落的规模将会减小，而群落将获得更好的稳定性。
\item 通过提高草本植物的吸水效率，灌木能够削弱种间竞争，使不同草本植物的平均生物量相互接近，并提升生物多样性的稳定性。
\item 在第 5.4 节所给出的更为严酷的干旱条件下，群落的规模将下降 67\% 以上，演化的波动幅度将增大 36\% 以上。但无论天气周期如何不同，这些指标随物种数目增加所呈现的趋势仍与第 5.2 节相似。
\item 我们比较了两类典型污染物的影响。结果表明，诸如酸雨之类的污染物更为危险，因为它会同时使群落规模下降 14.6\%、使波动增大 14.4\%。此外，栖息地减少 50\% 将使系统损失 22.9\% 的规模，但同时使波动降低 8.7\%。
\item 基于我们的模型与分析，我们的策略从两个主要方面加以表述：生态系统的保护（Preservation）与改造（Modification）。
\end{itemize}

最后，我们对模型的灵敏度与稳定性进行了充分分析。

\vspace{6pt}
\noindent\textbf{关键词：}生长模型；竞争模型；促进模型；四阶 Runge-Kutta 方法；TOPSIS

\newpage
\pagestyle{fancy}

%%====================== 正文 ======================
\section{引言}

\subsection{研究背景}

干旱环境的特点是降雨有限且多变，其资源供给呈脉冲式。以草本植物与灌木为主要覆被的草原，对水分短缺十分敏感。大量观测揭示：不同物种的数目，对植物群落如何适应那些以多变频率、在很宽的强度范围内发生的干旱，具有至关重要的影响。具体而言，在包含若干物种的群落中，植物对干旱条件的适应程度要优于仅含单一物种的群落。因此，这一由局地生物多样性所导致的现象具有重要的研究价值。

\subsection{问题的重述}

结合问题陈述中所给出的背景信息与限制条件，我们需要解决以下问题：

\begin{itemize}
\item 建立一个数学模型，用以描述植物群落中物种数目与干旱适应性之间的关系。特别地，我们需要预测当植物群落暴露于变化的天气周期时将会发生何种变化，并且我们的模型应能刻画这一过程中不同物种之间的相互作用。
\item 找出植物群落从局地生物多样性中获益所必需的最小物种数目，并预测当物种数目增加时将会发生什么。此外，我们还应通过改变植物群落中的物种类型来分析模型的灵敏度。
\item 讨论若干旱的发生具有波动的频率与更大的变异程度，物种数目对总体种群是否仍具有相同的影响。同时，我们还应把污染与栖息地缩减等其他因素纳入考虑。
\item 阐明为确保植物群落的长期存续应采取何种措施，并总结这一现象对更大范围环境的影响。
\end{itemize}

\subsection{本文的工作}

在本文中，我们逐层建立三个数学模型并回答六个问题，其结构如下所述。首先，我们建立干旱条件下单一物种的生长模型，用以刻画植物在随机降雨驱动下的生长与种子库更新；其次，在该模型基础上建立多物种的竞争模型，用以描述物种间对水资源的竞争及其对群落规模与稳定性的影响；最后，引入灌木植物对草本植物的促进模型，并进一步把干旱频率的波动、污染与栖息地缩减等因素纳入模拟。相应地，六个问题依次为：群落对可变天气周期的响应、群落获益所需的最小物种数目、物种数目继续增加时的演化趋势、干旱频率波动下的影响、污染与栖息地缩减等其他因素的影响，以及长期存续所需采取的措施。

\section{假设与合理性说明}

为简化问题，我们作出以下假设，并对每一条假设给出了合理性说明。

\begin{itemize}
\item \textbf{假设 1：}本文的研究对象主要是夏季型一年生植物，其生命周期由春季延续至秋季。这些植物的生物量在冬季将转化为种子库的生物量。

$\rightarrow$ \textbf{合理性说明：}根据文献 [1]，一年生植物具有生长迅速、繁殖能力强、生命周期短的特点。这类植物是干旱地区植被的重要组成部分，具有很高的研究价值。

\item \textbf{假设 2：}干旱地区的降雨是瞬时完成的，两次降雨之间不存在额外的水分供给。

$\rightarrow$ \textbf{合理性说明：}干旱地区每次降雨的持续时间，远短于我们所要研究的以年为尺度的时间跨度。与此同时，绝大多数干旱地区缺乏河流、湖泊等地表水资源。因此，我们可以认为降雨是唯一的水分供给途径。

\item \textbf{假设 3：}干旱地区的水分蒸发与土壤中的总水量成正比。

$\rightarrow$ \textbf{合理性说明：}干旱地区水分的蒸发作用全年都很显著。此外，干旱地区的土壤通常较为疏松，土壤中的水分均匀地暴露于空气之中，因此我们可以推断蒸发量与总水量成正比。

\item \textbf{假设 4：}一年生植物具有规律的开花结实期，且花期较短。

$\rightarrow$ \textbf{合理性说明：}一年生植物具有固定的生命周期与规律的花期。受干旱环境的影响，植物的花期一般较短，往往只有数日。这一时间跨度同样远短于我们所要研究的以年为尺度的时间跨度。

\item \textbf{假设 5：}植物的平均生物量（或平均质量）与根系的体积成正比。

$\rightarrow$ \textbf{合理性说明：}可以看出，植物的生物量越大，所需的资源就越多；而在干旱地区，植物的主要限制因素是水分短缺。考虑到根系是植物吸收水分的主要途径，假设植物的平均生物量与根系所覆盖的体积成正比是合理的。

\item \textbf{假设 6：}在实验区域内，植物均匀分布。

$\rightarrow$ \textbf{合理性说明：}为了减少对资源的竞争，植物倾向于在整个区域内均匀分布。

\item \textbf{假设 7：}在植物生长过程中（即萌发至开花之间），不考虑整株植物的死亡。

$\rightarrow$ \textbf{合理性说明：}在我们的模型中，环境限制与种间竞争被认为会影响植物的平均生物量。我们忽略上述因素可能导致整株植物死亡的可能性。
\end{itemize}

\section{符号说明}

本文所使用的主要符号罗列如下（同一符号在全文中的含义固定不变）。

\begin{itemize}
\item $r$：内禀增长率（inherent growth rate）；
\item $\bar m_i$：物种 $i$ 的平均生物量；
\item $K_i$：环境限制下物种 $i$ 的最大平均生物量；
\item $W$：资源可获得性，即土壤含水量；
\item $u(\cdot)$：单位生物量上的资源吸收速率；
\item $c$：由水转化为生物量增量的转化效率；
\item $m$：单位时间内植物的损失率；
\item $a$：吸水速率随资源可获得性变化的程度；
\item $u_{\max}$：吸水速率的上限；
\item $k$：控制吸收曲线变化趋势的参数；
\item $S_i$：物种 $i$ 的种子库生物量；
\item $\eta$：总生物量转化为种子的比例；
\item $\sigma$：种子在翌年萌发的比例；
\item $\alpha_{ij}$：单位生物量的植物 $i$ 所消耗的水资源，是植物 $j$ 的 $\alpha_{ij}$ 倍，其中植物 $j$ 位于其邻域；
\item $D_i$：植物 $i$ 根系伸展的平均直径；
\item $\beta_i$：相对水分利用能力；
\item $\bar A_i$：植物 $i$ 的平均根系覆盖面积；
\item $\bar A_i^{ov}$：植物 $i$ 与其他植物的平均重叠面积；
\item $V$：单株植物根系的体积；
\item $\lambda$：根系体积与生物量之间的比例因子；
\item $n_i$：植物（物种 $i$）的株数；
\item $u_{\max}'$：存在竞争时的吸水速率上限（限制项）；
\item $\gamma$：单位生物量灌木的提水能力；
\item $L$：根系的总长度；
\item $l$：单条根的平均长度；
\item $b$：根系系数，反映根系中的分枝数目；
\item $M^{(y)}$：第 $y$ 年植物物种的总重量；
\item $\bar M$：平均总重量；
\item $s^{2}$：总重量的逐年波动，即数据的方差。
\end{itemize}

\section{模型的建立}

\subsection{干旱环境下单一物种的生长模型}

本节建立仅有单一物种在干旱环境中的生长模型。此时，某一区域内只存在一个物种 $i$。我们定义 $\bar m_i$ 为物种 $i$ 的平均生物量，它将按照 Logistic 模型增长，即有
\[
\frac{d\bar m_i}{dt}=r\,\bar m_i\left(1-\frac{\bar m_i}{K_i}\right)
\]
其中 $r$ 为内禀增长率，$K_i$ 为环境限制下物种 $i$ 的最大平均生物量。

在干旱与半干旱地区，植物具有较为独特的生长规律，极端的气候条件与水资源约束对植物的正常生长造成严重影响 [2][3]。因此，有必要依据干旱地区的环境特征对上述 Logistic 方程加以修正。

\begin{itemize}
\item \textbf{自然增长率对干旱条件的响应}

根据关于干旱区植物生长的研究 [4]，植物的内禀增长率与植物自身的性状以及当地土地的土壤含水量有关，可表示为
\[
r=c\,u(W)-m
\]
其中 $r$ 为内禀增长率，$W$ 为资源可获得性（土壤含水量），$u(\cdot)$ 是单位生物量上的资源吸收速率，它是资源可获得性的函数，$c$ 表示植物把吸收的水分转化为生物量增量的转化效率，$m$ 是单位时间内植物的损失率，它由新陈代谢、自然死亡以及其他生物的破坏所引起。

特别地，此处的 $u(\cdot)$ 可写成如下形式
\[
u(W)=u_{\max}\,\frac{W^{a}}{W^{a}+k^{a}}
\]
其中 $a$ 表示吸水速率随资源可获得性变化的程度，$u_{\max}$ 控制吸水速率的上限，$k$ 控制吸收曲线的变化趋势。

\item \textbf{随机降雨模型}

正如假设 2 所述，干旱地区的降雨可视为一次瞬时的水分脉冲。因此，我们需要两个变量来确定一次具体的降雨。具体而言，$t_r$ 表示该次降雨在一年内的相对时刻，$R$ 表示该次降雨的总水量。

文献 [3] 给出了一个用于预测干旱地区一年内降雨脉冲出现时刻与降雨量的模型。首先，降雨的出现时刻是从参数为 $\alpha$ 与 $\beta$ 的 Beta 分布中随机抽取的，可写为 [5]
\[
f(t_r)=\frac{t_r^{\alpha-1}\,(1-t_r)^{\beta-1}}{\mathrm{B}(\alpha,\beta)},\qquad
\mathrm{B}(\alpha,\beta)=\frac{\Gamma(\alpha)\,\Gamma(\beta)}{\Gamma(\alpha+\beta)}
\]
其中 $\Gamma(\cdot)$ 为 Gamma 函数，$f(t_r)$ 为降雨时刻在一年内的概率密度。

降雨量同样被视为从对数正态分布中随机抽取的取值，其参数 $\mu$ 与 $\sigma^{2}$ 分别是降雨量自然对数的均值与方差。于是，$R$ 的概率密度可写为
\[
f(R)=\frac{1}{R\,\sigma\sqrt{2\pi}}\exp\left[-\frac{\left(\ln R-\mu\right)^{2}}{2\sigma^{2}}\right]
\]

\item \textbf{土壤含水量的变化规律}

根据假设 2，两次降雨之间不存在额外的水分供给。因此，土壤中的水分持续减少，其原因主要是植物的吸收与蒸发。特别地，如假设 3 所述，水分蒸发与总水量成正比。设比例因子为 $k_e$，则有
\[
\frac{dW}{dt}=-\,k_e\,W-\sum_{i}n_i\,\bar m_i\,u_i(W)
\]
与此同时，我们把降雨发生的时刻记为 $t_r$。根据假设 2，储存在土壤中的水分瞬时发生变化
\[
W\left(t_r^{+}\right)=W\left(t_r^{-}\right)+R
\]

\item \textbf{植物生命规律与种子库的演化}

根据假设 4，夏季型一年生植物具有规律的花期。在这一过程中，植物完成开花、交配与结实。一般而言，一年生植物在花期之后枯萎死亡，但其种子将以一定比例（记为 $\sigma$）在翌年生长。与此同时，总生物量中有一定比例转化为新的种子并补充进入种子库。设总生物量转化为种子的比例为 $\eta$，则可推得物种 $i$ 的种子库变化如下
\[
S_i^{(y+1)}=(1-\sigma)\,S_i^{(y)}+\eta\,n_i\,\bar m_i^{(y)}\left(t_f\right)
\]
其中 $t_f$ 表示花期的结束时刻。
\end{itemize}

综上，我们得到了物种 $i$ 在没有与其他物种竞争时的干旱环境演化模型。具体而言，其演化规则如下：

\begin{enumerate}
\item[(i)] 在生长过程中，平均生物量 $\bar m_i$ 与土壤含水量 $W$ 服从如下微分方程
\[
\begin{cases}
\dfrac{d\bar m_i}{dt}=r_i\,\bar m_i\left(1-\dfrac{\bar m_i}{K_i}\right)\\[1.0em]
\dfrac{dW}{dt}=-\,k_e\,W-n_i\,\bar m_i\,u_i(W)
\end{cases}
\]
\item[(ii)] 当降雨发生于 $t_r$ 时刻，土壤含水量 $W$ 将瞬时改变
\[
W\left(t_r^{+}\right)=W\left(t_r^{-}\right)+R
\]
\item[(iii)] 花期结束之后，植物种群产生种子。于是，总生物量与种子库由如下方程描述
\[
\begin{cases}
S_i\leftarrow S_i+\eta\,n_i\,\bar m_i\left(t_f\right)\\[0.6em]
M_i=n_i\,\bar m_i\left(t_f\right)
\end{cases}
\]
\item[(iv)] 花期结束之后，所有夏季型一年生植物都将死亡；在翌年春季，一部分种子按比例 $\sigma$ 萌发并开始生长，即
\[
n_i\,\bar m_i\big|_{t=0}=\sigma\,S_i^{(y)}
\]
\end{enumerate}

至此，单一物种在干旱环境中的生长模型（含种子库的逐年更新机制）已完整建立。该模型以随机降雨为外部驱动，以土壤含水量为资源状态变量，以平均生物量为种群状态变量，并借助花期的瞬时更新与年度间的种子萌发把年内生长与年际演化连接起来。模型中的内禀增长率并非常数，而是随土壤含水量变化的函数，这一点正是干旱环境区别于温和环境的关键所在：水分充足时增长率接近其上限 $c\,u_{\max}-m$，而水分匮乏时增长率迅速下降，甚至变为负值，从而使种群在干旱年份出现生物量的净损失。此外，由于降雨的时刻与水量均为随机变量，模型的解不再是确定性轨迹，而是一族依赖于降雨序列的样本轨道；这也意味着群落演化的“稳定性”必须以多条轨道上的统计量（均值与方差）来度量，而不能仅凭单条轨迹判断。
\subsection{多物种的竞争模型}

如前所述，我们已建立了单物种抵抗不规则气候变化的生长模型。本节把多个物种之间的竞争纳入考虑。我们先考察两个物种之间的竞争。在 Lotka-Volterra 模型的基础上，我们引入影响因子 [6]，在生长模型的基础上表达这种竞争
\[
\begin{cases}
\dfrac{d\bar m_1}{dt}=r_1\,\bar m_1\left(1-\dfrac{\bar m_1+\alpha_{12}\,\bar m_2}{K_1}\right)\\[1.0em]
\dfrac{d\bar m_2}{dt}=r_2\,\bar m_2\left(1-\dfrac{\bar m_2+\alpha_{21}\,\bar m_1}{K_2}\right)
\end{cases}
\]
其中 $\alpha_{12}$ 与 $\alpha_{21}$ 表示两个种群之间的竞争。考虑到在干旱环境中竞争的主要资源是水分，我们对 $\alpha_{ij}$ 作如下定义：

\begin{itemize}
\item $\alpha_{ij}$ 表示单位生物量的植物 $i$ 所消耗的水资源，是植物 $j$ 的 $\alpha_{ij}$ 倍，其中植物 $j$ 位于植物 $i$ 的邻域。
\end{itemize}

我们从根系生长的角度推导 $\alpha_{12}$ 与 $\alpha_{21}$ 的取值 [7]。假设植物 $i$ 根系伸展的平均直径为 $D_i$，则其平均半径为 $D_i/2$，于是根系覆盖面积为 $\pi D_i^{2}/4$。若植物 1 与植物 2 的根系覆盖范围发生重叠，则可推断这两株植物之间存在对水资源的竞争。

首先对水分利用能力参数作归一化处理，得到相对水分利用能力 $\beta_i$，它们满足
\[
\beta_i=\frac{\gamma_i}{\sum_{j=1}^{N}\gamma_j},\qquad \sum_{i=1}^{N}\beta_i=1
\]
其中 $\gamma_i$ 为物种 $i$ 的水分利用能力，$N$ 为该区域内物种的数目。

根据 $\alpha_{ij}$ 的定义，可以推得
\[
\alpha_{12}=\frac{\beta_1}{\beta_2}\cdot\frac{\bar A_1^{ov}}{\bar A_1},\qquad
\alpha_{21}=\frac{\beta_2}{\beta_1}\cdot\frac{\bar A_2^{ov}}{\bar A_2}
\]
其中 $\bar A_1^{ov}$ 表示植物 1 与其他植物的平均重叠面积，$\bar A_2^{ov}$ 表示植物 2 与其他植物的平均重叠面积，$\bar A_1=\pi D_1^{2}/4$ 为植物 1 的平均根系覆盖面积，$\bar A_2=\pi D_2^{2}/4$ 为植物 2 的平均根系覆盖面积。

至于根系覆盖面积 $\bar A_i$ 与平均重叠面积 $\bar A_i^{ov}$，我们可以依据实验区域内植物的数量与平均生物量对它们作出估计。

根据假设 5，我们可以由植物的平均生物量求出根系的体积
\[
V_i=\lambda\,\bar m_i
\]
其中 $\lambda$ 为比例因子，$V_i$ 为单株植物 $i$ 的根系体积，二者均可由植物的生长数据估计得到。具体而言，我们可以采集植物生长过程中的根长与生物量数据，据此确定该比例因子。

由根系体积，我们可以进一步得到根系覆盖面积
\[
A_i=\pi\left(\frac{3V_i}{2\pi}\right)^{2/3}
\]
该式由“把根系所占空间视为半球”的几何假设推得：半球体积为 $\frac{2}{3}\pi r^{3}$，故半径 $r=\left(\frac{3V}{2\pi}\right)^{1/3}$，覆盖面积为 $\pi r^{2}$。

如假设 7 所述，我们假定植物的株数（记为 $n_i$）在生长过程（萌发至开花之间）中保持不变。于是，基于假设 6，我们按下式估计各植物的平均重叠面积
\[
\bar A_i^{ov}=\max\left\{0,\ \frac{A_i\displaystyle\sum_{j\neq i}n_jA_j}{A_{\mathrm{total}}}\right\}
\]
其中 $A_{\mathrm{total}}$ 为实验区域的总面积。值得注意的是，若平均重叠面积的计算结果小于 0，则我们推断该植物与其他植物之间不存在竞争。

据此，我们整理出两个植物种群在一年生长期内相互竞争的微分方程模型，可描述为
\[
\begin{cases}
\dfrac{d\bar m_1}{dt}=r_1\,\bar m_1\left(1-\dfrac{\bar m_1+\alpha_{12}\,\bar m_2}{K_1}\right)\\[1.0em]
\dfrac{d\bar m_2}{dt}=r_2\,\bar m_2\left(1-\dfrac{\bar m_2+\alpha_{21}\,\bar m_1}{K_2}\right)\\[1.0em]
\dfrac{dW}{dt}=-\,k_e\,W-\left(n_1\,\bar m_1\,u_1(W)+n_2\,\bar m_2\,u_2(W)\right)
\end{cases}
\]
其中各参数可由如下规则确定：
\[
\begin{cases}
r_i=c_i\,u_i(W)-m,\qquad u_i(W)=u_{\max}\,\dfrac{W^{a}}{W^{a}+k^{a}}\\[0.8em]
\alpha_{12}=\dfrac{\beta_1}{\beta_2}\cdot\dfrac{\bar A_1^{ov}}{\bar A_1},\qquad
\alpha_{21}=\dfrac{\beta_2}{\beta_1}\cdot\dfrac{\bar A_2^{ov}}{\bar A_2}\\[0.8em]
\alpha_{12}\,\alpha_{21}<1
\end{cases}
\]
最后一条不等式保证了两个种群能够共存于同一平衡态，而不会出现一方完全排斥另一方的情形。

进一步推广，我们得到 $n$ 个植物种群在一年生长期内相互竞争的微分方程模型，可描述为
\[
\begin{cases}
\dfrac{d\bar m_i}{dt}=r_i\,\bar m_i\left(1-\dfrac{\bar m_i+\displaystyle\sum_{j\neq i}\alpha_{ij}\,\bar m_j}{K_i}\right),\qquad i=1,2,\dots,n\\[1.4em]
\dfrac{dW}{dt}=-\,k_e\,W-\displaystyle\sum_{i=1}^{n}n_i\,\bar m_i\,u_i(W)
\end{cases}
\]
其中各参数可由如下规则确定：
\[
\begin{cases}
r_i=c_i\,u_i(W)-m,\qquad u_i(W)=u_{\max}\,\dfrac{W^{a}}{W^{a}+k^{a}}\\[0.8em]
\alpha_{ij}=\dfrac{\beta_i}{\beta_j}\cdot\dfrac{\bar A_i^{ov}}{\bar A_i}\\[0.8em]
K_i\ \text{由物种自身的性状决定，}\ \beta_i\ \text{由归一化得到}
\end{cases}
\]
在以上微分方程的基础上，把降雨的瞬时更新、花期后的种子产生以及年度间的种子萌发一并考虑，即可得到最终的模型。

值得强调的是，上述竞争模型的结构具有明确的生态学含义：括号中的 $\bar m_i+\sum_{j\neq i}\alpha_{ij}\bar m_j$ 称为“有效生物量”，它把其他物种的生物量按其对水分的相对消耗能力折算为本物种的等价生物量。由于 $\alpha_{ij}$ 与根系重叠面积成正比，当两个物种的根系在空间上几乎不重叠（例如深根与浅根物种）时，$\alpha_{ij}$ 接近 0，物种之间的竞争被极大地削弱；反之，若两个物种的根系在浅层土壤中高度重叠且水分利用能力相近，则 $\alpha_{ij}$ 接近 1，模型退化为两个几乎完全对称竞争的种群。这一机制使得模型能够自然地反映“生态位分化削弱竞争”这一生态学共识，而无需额外引入人为的竞争强度参数。

\subsection{灌木植物对草本植物的促进模型}

在第 4.2 节中，我们从水分根系竞争的角度建立了干旱周期内植物的种间竞争模型。该模型适用于所有植物。然而，面对干旱胁迫，灌木植物在促进草本植物生长方面具有独特的作用 [8]。这种促进作用主要体现在灌木能够改善草本植物的存活与生长。考虑到引入灌木并不能直接提高土壤含水量 $W$，我们推断灌木能够改善吸水速率的限制因子（记为 $u_{\max}'$）。根据关于吸水速率的表达式可见，$u$ 受 $a$ 与 $u_{\max}$ 的影响。由于 $a$ 表示吸水速率随资源可获得性变化的程度，它是植物的一种固有属性，我们认为 $a$ 不会因灌木的引入而改变。因此，我们推测灌木是通过影响 $u_{\max}$ 来起作用的。在第 4.2 节竞争模型的基础上，我们进一步把参数 $u_{\max}$ 修正为 $u_{\max}'$。

根据 Richards 团队 [9] 与 Prieto 团队 [10] 的研究，我们发现灌木的根系能够把水分从深层土壤提升到浅层土壤。这一过程可以改善浅层土壤的环境，有利于其他植物吸收水分。考虑到该限制项取值越小、植物的吸水速率越高，我们推断
\[
u_{\max,i}'=\frac{u_{\max,i}}{1+\gamma\displaystyle\sum_{s}n_s\,\bar m_s}
\]
其中 $\gamma$ 为比例因子，它反映单位生物量灌木的提水能力，求和遍历植株 $i$ 邻域内的所有灌木。

对该式的分析表明：当植株 $i$ 周围没有灌木时，$u_{\max,i}'$ 等于 $u_{\max,i}$；当存在灌木时，灌木的生物量越大，$u_{\max,i}'$ 越小，因此若植株 $i$ 周围有灌木，其吸水速率 $u_i$ 将会提高。

与此同时，根据 Butterfield 团队 [11] 的研究，灌木对土壤水分的促进作用会随着干旱的加剧而减弱。考虑到土壤含水量是衡量干旱程度的判据之一，我们进一步把上式修正为
\[
u_{\max,i}'=\frac{u_{\max,i}}{1+\gamma\,W\displaystyle\sum_{s}n_s\,\bar m_s}
\]
其中 $W$ 为资源可获得性（土壤含水量）。随着土壤含水量的增大，$u_{\max,i}'$ 单调减小，于是 $u_i$ 增大，这意味着植物以更高的速率吸收水分；而在土壤极为干燥时，灌木的提水效果被显著放大，模型由此再现了“干旱胁迫越强、促进作用越显著”的经验规律。

由于灌木对草本植物的促进机制复杂且具有不确定性，目前尚无较为一致的研究结论 [8]。因此，我们依据对具体情形的分析建立上述模型，并将在模拟中给出对参数的自有估计。

\subsection{关键参数的合理估计}

本节对前述模型中的参数作出合理估计。

\begin{itemize}
\item \textbf{随机降雨模型中的统计参数}

根据第 4.1 节所建立的随机降雨模型，一年内的平均降雨量可表示为 $\exp\left(\mu+\sigma^{2}/2\right)$。由于 $\sigma$ 与 $\mu$ 均未知，我们需要确定均值 $\mu$ 与变异程度 $\sigma$。

我们查阅了不同干旱程度地区的年降水量。我们发现，重度干旱地区的年平均降水量仅为 32.5 mm [12]，例如塔克拉玛干沙漠南缘；中度干旱地区的年平均降水量约为 113.2 mm [13]，例如中国河西走廊中部的原始荒漠；而半干旱地区的这一数据约为 360 mm [14]，例如中国内蒙古乌审旗的沙地柏自然保护区。

考虑到我们的实验区域为 10000 $\mathrm{m^{2}}$，则对应于上述不同干旱程度，年平均降水质量分别为 325 kg、1132 kg 与 3600 kg。依据统计学中的 $3\sigma$ 原则的常用近似，有 95\% 的概率使降雨量落在均值两侧两倍标准差的范围之内，即
\[
\begin{cases}
\exp\left(\mu+\dfrac{\sigma^{2}}{2}\right)=\bar R\\[1.0em]
\exp\left(\mu+\dfrac{\sigma^{2}}{2}\right)\sqrt{e^{\sigma^{2}}-1}=\dfrac{1}{4}\,\bar R
\end{cases}
\]
其中 $\bar R$ 为该地区的年平均降雨量，第二式左端即对数正态分布的标准差。

上述关于 $\mu$ 与 $\sigma$ 的方程是超定的（三组数据对应六个方程，而未知量只有两个）。依据正规方程，我们基于最小二乘法求得最优参数值
\[
\begin{bmatrix}\mu\\[0.3em] \sigma^{2}\end{bmatrix}=\left(A^{\mathrm{T}}A\right)^{-1}A^{\mathrm{T}}B
\]
其中 $A$ 为上述方程中关于 $\mu$ 与 $\sigma^{2}$ 的系数矩阵，$B$ 为由各观测数据构成的右端向量。

最终求得 $\mu\approx 6.97$，$\sigma\approx 0.25$。这一取值意味着单次降雨量的典型规模在数百千克量级，且其变异程度约为均值的四分之一，与干旱区“稀遇而集中的脉冲式降水”这一定性认识相符。

\item \textbf{根系体积与植物生物量之间的比例因子 $\lambda$}

在干旱环境中，植物根系会尽可能覆盖更大的面积。根据关于根系拓扑结构的讨论 [13]，一个理想的模型是二叉分枝。

为便于进一步讨论，我们引入根与根连接长度的概念。红砂（\textit{Reaumuria soongorica}，RS）与白刺（\textit{Nitraria tangutorum}，NT）是荒漠地区两个常见的种群。根据文献 [12]，RS 的根系总长度平均为 3185.5 cm，平均根连接长度为 29.8 cm；而 NT 的根系总长度平均为 4802.5 cm，平均根连接长度为 49.0 cm。

通过分析二叉分枝的特征，我们有
\[
L=b\,l
\]
其中 $L$ 为根系总长度，$l$ 为单条根的平均长度，$b$ 为根系系数，它反映根系中分枝的数目。

接下来，我们把根系总长度与平均根连接长度代入上式，得到
\[
b_{\mathrm{RS}}=\frac{3185.5}{29.8}\approx 106.9,\qquad
b_{\mathrm{NT}}=\frac{4802.5}{49.0}\approx 98.0
\]
随后，我们考虑植物根系会以球形向四周伸展以吸收水分，于是按下式估计根系的直径
\[
D=l\sqrt{b}
\]
该式的含义是：在二叉分枝的结构下，单条根的长度 $l$ 与分枝数目的平方根共同决定了根系在空间中的伸展尺度。代入数据得 $D_{\mathrm{RS}}\approx 3.08$ m，$D_{\mathrm{NT}}\approx 4.85$ m，与荒漠灌木根系水平伸展数米的观测结果一致。

我们把根系简化为半个球体，从而也可得到根系的体积
\[
V=\frac{2}{3}\pi\left(\frac{D}{2}\right)^{3}
\]
再把两个种群单株植物的质量纳入考虑，我们由根系体积与生物量之间的比例关系确定 $\lambda_{\mathrm{RS}}$ 与 $\lambda_{\mathrm{NT}}$。由于这两个种群在干旱地区具有代表性，我们依据二者取值的平均来估计通用的比例因子 $\lambda$，即
\[
\lambda=\frac{\lambda_{\mathrm{RS}}+\lambda_{\mathrm{NT}}}{2}\approx 2.7\times 10^{7}\ \mathrm{cm^{3}\cdot kg^{-1}}
\]
由于 $\lambda$ 是根系体积与植物生物量之间的比例因子，它反映的是一种普遍规律。因此，我们认为不同植物的这一比例因子近似相等。在模拟中，我们用 $\lambda$ 统一替代不同物种的比例因子 $\lambda_i$。

\item \textbf{其他参数的估计}

除上述两个参数之外，我们还需要确定 $c$、$a$ 与 $K_i$。

如前所述（见第 4.1 节），$c$ 表示植物把吸收的水分转化为生物量增量的转化效率，$a$ 表示吸水速率随资源可获得性变化的程度。根据文献 [4]，$c$ 的合理取值范围为 $c\in[0.2,\,0.6]$，而 $a$ 的取值范围为 $a\in[1,\,4]$。在考虑不同的植物物种时，我们在上述范围内取值，以确保最终结果的合理性。

根据第 4.1 节的前期分析，$K_i$ 是环境限制下物种 $i$ 的最大平均生物量。显然，不同物种的 $K_i$ 应当不相等。然而，对于本文所研究的夏季型一年生植物，$K_i$ 一般小于 0.5 kg [7]。因此，把其取值设定在 0.3 kg 附近是合理的。在后续模拟中，除灵敏度分析外，我们取各物种的 $K_i$ 相同。
\end{itemize}

\section{模型的计算与结果分析}

\subsection{模型的求解算法}

考虑到四阶 Runge-Kutta 算法具有较高的精度与较强的稳定性，我们采用该方法计算前述微分方程组的数值解。首先，我们把微分方程组整理为 $\dfrac{d\mathbf{y}}{dt}=\mathbf{f}(t,\mathbf{y})$ 的形式，其中 $\mathbf{y}=\left(\bar m_1,\dots,\bar m_n,W\right)^{\mathrm{T}}$。然后，我们构造如下差分迭代格式以计算一阶微分方程
\[
\begin{cases}
\mathbf{y}_{k+1}=\mathbf{y}_k+\dfrac{h}{6}\left(\mathbf{K}_1+2\mathbf{K}_2+2\mathbf{K}_3+\mathbf{K}_4\right)\\[0.8em]
\mathbf{K}_1=\mathbf{f}\left(t_k,\mathbf{y}_k\right),\qquad
\mathbf{K}_2=\mathbf{f}\left(t_k+\dfrac{h}{2},\mathbf{y}_k+\dfrac{h}{2}\mathbf{K}_1\right)\\[0.8em]
\mathbf{K}_3=\mathbf{f}\left(t_k+\dfrac{h}{2},\mathbf{y}_k+\dfrac{h}{2}\mathbf{K}_2\right),\qquad
\mathbf{K}_4=\mathbf{f}\left(t_k+h,\mathbf{y}_k+h\mathbf{K}_3\right)
\end{cases}
\]
其中 $h$ 为计算过程的时间步长。

计算流程如下所述：

\begin{itemize}
\item \textbf{输入：}（1）各物种的性状参数与群落的初始条件；（2）模拟的时间跨度；（3）基于 Beta 分布模型生成的降雨时刻；（4）基于对数正态分布模型生成的每次降雨的水量。
\item \textbf{循环：}当模拟时间尚未达到设定跨度时，反复执行以下步骤：
\begin{enumerate}
\item 若处于一年的开始，则种子库中特定比例的种子萌发；
\item 使用 Runge-Kutta 方法计算该时刻各物种生物量的增量 $\Delta \bar m_i$ 以及地下总储水量 $W$；
\item 若遇到降雨，则用相应的水量 $R$ 增加储水量 $W$；
\item 若该年的花期已经结束，则把生物量 $\bar m_i$ 的 $\eta$ 比例转化为种子库的增量，随后所有一年生植物死亡；
\item 记录群落的演化状态。
\end{enumerate}
\end{itemize}

\subsection{多植物群落的效益评价}

对于一个植物群落而言，其面对不规则天气条件的表现可从两个方面评价。其一是所有植物的重量，它应当尽可能高；其二是应当避免极端波动，以保证生态系统的稳定性。总之，一个理想的生态系统应当在规模与稳定性之间取得平衡。

群落在第 $y$ 年的规模可建模为植物物种的总重量（记为 $M^{(y)}$），可写成
\[
M^{(y)}=\sum_{i=1}^{N^{(y)}}n_i\,\bar m_i^{(y)}
\]
其中 $\bar m_i^{(y)}$ 为植物 $i$ 在第 $y$ 年的平均生物量，$N^{(y)}$ 为该年植物物种的总数。为简化问题的分析，我们把时间限制在 $Y$ 年的区间内。

考虑每一年植物的总重量，我们可以计算该区间内的平均总重量（记为 $\bar M$），即
\[
\bar M=\frac{1}{Y}\sum_{y=1}^{Y}M^{(y)}
\]
总重量的逐年波动可视为数据的方差（记为 $s^{2}$），其表达式为
\[
s^{2}=\frac{1}{Y}\sum_{y=1}^{Y}\left(M^{(y)}-\bar M\right)^{2}
\]
若生物多样性确实为植物群落带来了益处，则它应当体现为更高的规模或更低的波动。为评价这两方面之间的平衡，我们采用 TOPSIS 方法。TOPSIS（逼近理想解的排序技术，Technique for Order Preference by Similarity to an Ideal Solution）是多目标决策分析中一种常用的方法。因此，我们采用这一模型评价植物群落演化的表现。若评价表现优于仅有一个植物物种的情形，则说明该植物群落已从生物多样性中获益。由此，我们能够确定需要多少个物种才能使群落获益。

与此同时，该模型也有助于揭示物种数目增加时将会发生什么。我们所需做的只是扩大物种数目并将其作为模拟程序的输入。

区域内的物种越多，植物的总重量越低，而群落面对不规则天气条件所经历的波动越小。这一现象可由前述微分方程模型直观地解释：随着物种数目的增加，植物群落的生长将因对有限水资源的更激烈竞争而承受更多的“摩擦”。在这一竞争激烈的群落中，任何一类物种都受到抑制，难以在短时间内增长到极大的规模。因此，不同物种之间的相对比例得以维持，从而使植物群落具有更高的稳定性。

由 TOPSIS 导出的效益评价指标表明，当物种数目多于八个时，该指标将大于只有一个物种时的指标。因此，需要八种以上的物种，群落才能在规模与稳定性两方面获得整体效益。

\subsection{群落中不同物种类型的影响}

本节给出不同物种类型对植物群落演化的影响。基于第 4.2 节的竞争模型，我们发现植物群落中不同物种之间的 $\beta_i$ 差异越大，总生物量越大，但其稳定性越低。我们进一步依据第 4.3 节的模型考虑上述环境中灌木植物对其他植物的贡献。具体而言，我们模拟两类植物群落并比较它们的差异。

在第一种情形中，群落包含三种不同的草本植物与一种灌木。如第 4.3 节所述，灌木能够通过提高草本植物的吸水速率来促进其生长。因此，除物种之间的竞争之外，还存在合作关系。在第二种情形中，植物群落仅由三种草本植物构成，不含灌木。换言之，物种之间只有竞争。我们把这二者置于相同的天气条件之下，模拟结果显示：

根据对演化数据的分析，第一个群落（即含灌木者）表现出相对较低的平均总重量 $\bar M$，以及比第二个群落（不含灌木者）小得多的方差 $s^{2}$。这一差异表明，由于根系的竞争，灌木的加入导致群落规模相对降低；然而，面对不规则的天气条件，该植物群落具有高得多的稳定性。更重要的是，灌木能够削弱不同草本植物之间的竞争，这有利于生物多样性的维持。

\subsection{干旱发生频率波动的影响}

干旱天气条件通常表现为每次降雨量更少、降雨频次更低。为模拟各种干旱频次与发生情形对植物群落的影响，一种可行的做法是相应地改变各次降雨的时间分布与水量。当干旱发生时，降雨以相同频次出现的概率将会降低（甚至可能出现整年无雨的情况！）。这需要一个新随机变量 $N_{\mathrm{rain}}$ 来控制降雨次数。此外，每次降雨的水量也将更低且波动更大，这意味着更小的 $\mu$ 与更大的 $\sigma$。不同情形下的降雨分布随之改变。

在模拟中，我们假设 $N_{\mathrm{rain}}$ 服从参数为 $\lambda_p$ 的 Poisson 分布。与此同时，$\mu$ 与 $\sigma$ 也被改变，以构造更为严酷的干旱条件。

面对严酷的干旱条件，总体水资源的可获得性将进一步受限。因此，植物物种的演化不仅会带来更低的规模，也会带来更低的稳定性，这是可以理解的。此外，我们的模型表明，随着物种数目的增长，植物群落的规模与稳定性遵循（第 5.2 节所述的）相似趋势。具体而言，当类型数分别为 4、6、8、10、12 与 24 时，平均总重量分别为 209.06、184.75、162.48、153.6、147.22 与 117.14，而相应的方差分别为 26346、15428、13606、12569、7277.6 与 3986.1。可见规模随物种数目单调下降，波动亦单调下降，二者的变化方向与常态干旱下的结论完全一致；但与常态情形相比，同等物种数目下的规模整体更低、波动整体更大，这正对应摘要中“规模下降 67\% 以上、波动增大 36\% 以上”的定量结论。

\subsection{其他因素对植物群落的影响}

\begin{itemize}
\item \textbf{不同类型污染的影响}

一般而言，污染物会改变植物的新陈代谢，使其易受病害或虫害侵袭。本节研究两类典型污染物对植物群落的影响，即光化学烟雾与酸雨。它们对植物的具体损害如下 [15]：

\begin{itemize}
\item 光化学烟雾源于工业活动排放的氮氧化物与挥发性有机物（VOCs）之间的化学反应，其产物为地面臭氧与过氧乙酰硝酸酯。暴露于这些污染物中的植物更为脆弱，存活概率更低。因此，光化学烟雾的影响可建模为植物损失率的上升（记为 $m$ 的增大）。
\item 酸雨源于硫与二氧化氮同大气中其他物质的化学反应。直接暴露于酸雨会使植物更难进行光合作用。因此，酸雨的影响可视为转化效率的下降（记为 $c$ 的减小）。
\end{itemize}

为模拟污染对植物群落的影响，应对植物的损失率与转化效率作相应修正。含污染物的植物群落生长模拟结果显示：与非污染情形相比，暴露于污染之中的植物群落其规模更低。此外还值得注意的是，不同类型的污染物以不同方式影响系统的稳定性。对于提高生物量损失率的污染物（例如光化学烟雾），它并未对稳定性施加过多的负面影响；然而，当植物暴露于降低转化效率的酸雨时，情形则不同——此时植物群落的稳定性将显著低于非污染情形。这一现象可由水分利用在干旱地区的重要性加以说明：转化效率的下降直接削弱了植物把稀缺水分转化为生物量的能力，使群落对降雨波动的缓冲能力被系统性削弱。

\item \textbf{栖息地缩减的影响}

类似地，我们的模型也可用于探讨栖息地缩减对物种生长的影响。模拟结果表明，栖息地越小，规模越小。由于可供植物物种生长的空间更少，不同植物根系之间将出现更多的交叠，从而加剧竞争。
\end{itemize}

\subsection{长期存续所需的措施}

基于模拟结果，我们得到一系列确保植物群落长期存续的策略。我们的想法从两个方面表述，即生态系统的保护与改造。

\begin{itemize}
\item \textbf{保护（Preservation）}

植物群落的生态系统已在干旱地区存续了数百万年，这表明其面对严酷条件具有内在的稳健性。因此，应当采取措施保护群落的自然性。具体策略如下：

\begin{enumerate}
\item 保护植物群落的生物多样性。由第 4.2 节与第 4.3 节的模型模拟，我们发现多个植物物种之间的竞争与促进共同保障了系统面对严酷自然条件时的稳定性。
\item 保护干旱地区的水资源。若水资源被过度开采，将加剧水分短缺，从而降低植物群落的规模与稳定性。
\item 防治污染。由第 5.5 节的分析，污染物（尤其是酸雨一类）已被证明会损害植物群落的规模与稳定性。
\item 维护植物的栖息地。足够大的栖息地可为个体留出充足的生存空间，有助于避免物种之间的激烈竞争。
\end{enumerate}

\item \textbf{改造（Modifications）}

尽管天然草群落具有上述种种优势，它们面对严酷的干旱条件时仍然脆弱。通过改造可以提高植物群落的存续能力。具体策略如下：

\begin{enumerate}
\item 引入新物种。我们发现，诸如灌木一类的物种能够为其他物种带来促进作用，并为生态系统增添更多稳定性，这与文献的结论一致。因此，可以引入一些适宜的物种。
\item 建设灌溉设施。某些灌溉设施（例如伊朗的坎儿井）能够以较低的蒸发量收集并保存降雨形成的水分。因此，通过人工设施提高土壤水分的可获得性同样值得考虑。
\end{enumerate}
\end{itemize}

\section{灵敏度分析与稳定性分析}

\subsection{灵敏度分析}

在第 4 节中，我们建立了三个模型，其中包含许多参数。在第 4.4 节中，我们对大多数参数作出了合理估计。然而，仍有部分取值难以通过计算估计，我们只能参阅文献。本节对来自文献的参数进行灵敏度分析。

首先，$K_i$ 是物种 $i$ 的最大平均生物量，它反映环境对植物生长的限制。考虑到我们所研究的主要是干旱环境中的一年生植物，不同物种的最大平均生物量不应相差很大。因此，在此前的模拟中，我们假设不同物种具有相同的 $K_i$。通过调整 $K_i$ 的取值，我们得到灵敏度分析曲线（以三个物种的竞争为例），分别考察总生物量以及物种 1、物种 2、物种 3 的平均生物量随 $K_i$ 的变化。结果表明，当 $K_i$ 在其合理范围内变动时，总生物量与各物种生物量的定量结果虽有偏移，但随物种数目变化的单调趋势以及群落规模与稳定性之间的权衡关系保持不变。

随后，我们分析 $k$（控制吸收曲线变化趋势）与 $\eta$（植物生物量转化为种子的比例）的灵敏度。对 $k$ 而言，其取值决定吸水速率达到上限的速度：$k$ 越小，植物在低含水量下即可接近最大吸水速率，群落对轻度干旱的耐受性更强；$k$ 越大，吸水速率对含水量越敏感，群落规模的整体水平下降但波动格局不变。对 $\eta$ 而言，其取值决定当年生物量中有多大比例被“储存”为翌年的种子：$\eta$ 越大，种子库对年际干旱的缓冲作用越强，群落的长期规模越稳定；$\eta$ 过小则种子库难以补偿死亡年份的损失，群落可能走向退化。

可以看出，我们的模型在各种参数的灵敏度分析中表现良好。随着植物逐代繁殖，参数取值的波动对结果的影响逐渐减小。这表明我们的模型具有良好的长期稳定性。

\subsection{稳定性分析}

\begin{itemize}
\item \textbf{初始值对结果的影响}

我们选取模型的两个初始值，即初始平均生物量 $\bar m_i(0)$ 与初始株数 $n_i(0)$。我们对这两个取值作随机扰动，扰动变量服从均匀分布，其均值为 0，方差为
\[
\mathrm{Var}\left(\xi\right)=\frac{\left(2\delta\right)^{2}}{12}=\frac{\delta^{2}}{3}
\]
其中 $\delta$ 为扰动的幅度。

接下来，我们观察扰动对长期结果的影响。此处植物的总生物量取为一年生植物在死亡之前的最大生物量。可以看到，模型的长期稳定性是令人满意的：不同随机扰动下的轨道在若干年之后趋于同一统计意义上的水平，扰动所造成的影响随时间的推移而被繁殖过程逐渐“平均掉”。

\item \textbf{参数与结果稳定性之间的关系}

当我们讨论微分方程平衡点的稳定性时，通常先通过计算偏导数确定可能的平衡点，然后研究系数矩阵的特征方程并分析相轨线，最终确定方程在合理平衡点处稳定的条件 [6]。

基于第 4.2 节所建立的竞争模型，我们考虑两个物种之间的竞争。此时我们只考虑一年之内的生长，即忽略开花、结实与死亡的过程。由该节的方程，可以得到两个物种可能的平衡点
\[
\bar m_1^{*}=\frac{K_1-\alpha_{12}K_2}{1-\alpha_{12}\alpha_{21}},\qquad
\bar m_2^{*}=\frac{K_2-\alpha_{21}K_1}{1-\alpha_{12}\alpha_{21}}
\]
若希望两个物种的平均生物量在 $\left(\bar m_1^{*},\bar m_2^{*}\right)$ 处达到共存，我们可以用上述方法得到对参数取值的要求
\[
K_1>\alpha_{12}K_2,\qquad K_2>\alpha_{21}K_1,\qquad \alpha_{12}\,\alpha_{21}<1
\]
根据定义与平均重叠面积的估计式，平均根系覆盖面积 $\bar A_i$ 一般显著大于平均重叠面积 $\bar A_i^{ov}$。与此同时，对于耐旱的一年生植物，其相对吸水能力 $\beta_i$ 并无明显差异，可描述为
\[
\beta_1\approx\beta_2\ \Longrightarrow\ \alpha_{12}\approx\frac{\bar A_1^{ov}}{\bar A_1}\ll 1,\qquad
\alpha_{21}\approx\frac{\bar A_2^{ov}}{\bar A_2}\ll 1
\]
因此，上述参数条件在我们的模型中是合理的，且易于满足。
\end{itemize}

\section{模型的优缺点}

\subsection{优点}

\begin{itemize}
\item 我们的模型充分考虑了年生植物物种演化的复杂性，在 Logistic 模型的基础上采用微分方程的形式；同时，植物繁殖与死亡所引起的植物群落的瞬时跃变也被纳入其中。在理论上，我们的模型能够预测群落任意给定时刻的状态。
\item 我们在模型中借助自然增长率与土壤含水量之间的关系，考虑了水分可获得性对植物群落的影响。因此，我们的模型能够模拟植物群落在各种干旱条件下的演化。
\item 不同物种之间的促进与竞争，是依据已有的植物学研究与 Lotka-Volterra 模型建模的。因此，我们的模型能够模拟物种间相互作用各不相同的多种群落的演化。
\item 使用四阶 Runge-Kutta 算法计算微分方程的数值解，使我们的模型兼具计算精度与较强的稳定性。
\item 我们采用 TOPSIS 评价不同植物群落的演化。该方法兼具简洁性与可理解性，能够就规模与稳定性两方面高效地度量某一植物群落的整体表现。
\end{itemize}

\subsection{缺点}

\begin{itemize}
\item 在我们的模型中，土壤水分仅来自降雨，且仅因蒸发与物种的消耗而减少；而地下水的实际演化过程要复杂得多，地下河等因素不可忽略。
\item 各类干旱是以降雨频率与降水量的不同来定义的；然而干旱在许多其他方面也存在差异，例如局地温度、风环流等。因此，应当引入更多的维度以更好地刻画干旱。
\item 尽管模型中的多个参数保证了模型的灵活性，但其中一些参数难以精确计算，我们只能依据所参阅的文献给出恰当估计。不过，我们对模型中的关键参数进行了严格的推导与分析。
\end{itemize}

\section*{参考文献}

\begin{enumerate}
\item[{[1]}] Floristic composition, distribution, and resource types of desert annual plants in Xinjiang. DU Jiaqian [J]. Arid Zone Research, 2022, 39(01): 185-209.
\item[{[2]}] Leopold, A. C. (1964). Plant growth and development. Plant growth and development.
\item[{[3]}] Seleiman, M. F., Al-Suhaibani, N., Ali, N., Akmal, M., Alotaibi, M., Refay, Y., ... \& Battaglia, M. L. (2021). Drought stress impacts on plants and different approaches to alleviate its adverse effects. Plants, 10(2), 259.
\item[{[4]}] Chesson, P., Gebauer, R. L., Schwinning, S., Huntly, N., Wiegand, K., Ernest, M. S., ... \& Weltzin, J. F. (2004). Resource pulses, species interactions, and diversity maintenance in arid and semi-arid environments. Oecologia, 141, 236-253.
\item[{[5]}] Johnson NL, Kotz S, Balakrishnan N (1995) Continuous univariate distributions, vol 2, 2nd edn. Wiley, New York.
\item[{[6]}] JIANG Qi-yuan (2011). Mathmatical model. Higher Education Press.
\item[{[7]}] Fan, L., \& Zizhen, L. (2003). Interspecific competition model and application of plants in sandy area. Acta Botanica Boreali-occidentalia Sinica, 23(1), 138-140.
\item[{[8]}] LUO, G. S. C. T. X., \& Zhang, E. L. L. (2022). Advances in the study of shrub facilitation on herbs in arid and semi-arid regions. Chinese Journal of Plant Ecology, 46(11): 1321-1333.
\item[{[9]}] Richards JH, Caldwell MM (1987). Hydraulic lift: substantial nocturnal water transport between soil layers by Artemisia tridentata roots. Oecologia, 73, 486-489.
\item[{[10]}] Prieto I, Padilla FM, Armas C, Pugnaire FI (2011). The role of hydraulic lift on seedling establishment under a nurse plant species in a semi-arid environment. Perspectives in Plant Ecology Evolution and Systematics, 13, 181-187.
\item[{[11]}] Butterfield BJ, Bradford JB, Armas C, Prieto I, Pugnaire FI (2016). Does the stress-gradient hypothesis hold water? Disentangling spatial and temporal variation in plant effects on soil moisture in dryland systems. Functional Ecology, 30, 10-19.
\item[{[12]}] GUO, J. H., ZENG, F. J., LI, C. J., \& Zhang, B. (2014). Root architecture and ecological adaptation strategies in three shelterbelt plant species in the southern Taklimakan Desert. Chinese Journal of Plant Ecology, 38(1), 36.
\item[{[13]}] Li-shan, S. H. A. N., Yi, L. I., Wei, R. E. N., Shi-ping, S. U., Qiu-lian, D. O. N. G., \& Dong-mei, G. E. N. G. (2013). Root architecture of two desert plants in central Hexi Corridor of Northwest China. Yingyong Shengtai Xuebao, 24(1).
\item[{[14]}] He Weiming. Distribution Characteristic of Root Area of Sabina Vulgaris under Different Habitats. Scientia Silvae Sinicae, 2000, 36(5): 17-21.
\item[{[15]}] "Greentumble." [online]. Available: https://greentumble.com/effect-of-pollution-on-plants/
\end{enumerate}

\end{document}
